
\documentclass[a4paper,11pt]{article}
\usepackage[a4paper, margin=8em]{geometry}

% usa i pacchetti per la scrittura in italiano
\usepackage[french,italian]{babel}
\usepackage[T1]{fontenc}
\usepackage[utf8]{inputenc}
\frenchspacing 

% usa i pacchetti per la formattazione matematica
\usepackage{amsmath, amssymb, amsthm, amsfonts}

% usa altri pacchetti
\usepackage{gensymb}
\usepackage{hyperref}
\usepackage{standalone}

% imposta il titolo
\title{Appunti Comunicazioni Numeriche}
\author{Luca Seggiani}
\date{2026}

% imposta lo stile
% usa helvetica
\usepackage[scaled]{helvet}
% usa palatino
\usepackage{palatino}
% usa un font monospazio guardabile
\usepackage{lmodern}

\renewcommand{\rmdefault}{ppl}
\renewcommand{\sfdefault}{phv}
\renewcommand{\ttdefault}{lmtt}

% disponi il titolo
\makeatletter
\renewcommand{\maketitle} {
	\begin{center} 
		\begin{minipage}[t]{.8\textwidth}
			\textsf{\huge\bfseries \@title} 
		\end{minipage}%
		\begin{minipage}[t]{.2\textwidth}
			\raggedleft \vspace{-1.65em}
			\textsf{\small \@author} \vfill
			\textsf{\small \@date}
		\end{minipage}
		\par
	\end{center}

	\thispagestyle{empty}
	\pagestyle{fancy}
}
\makeatother

% disponi teoremi
\usepackage{tcolorbox}
\newtcolorbox[auto counter, number within=section]{theorem}[2][]{%
	colback=blue!10, 
	colframe=blue!40!black, 
	sharp corners=northwest,
	fonttitle=\sffamily\bfseries, 
	title=~\thetcbcounter: #2, 
	#1
}

% disponi definizioni
\newtcolorbox[auto counter, number within=section]{definition}[2][]{%
	colback=red!10,
	colframe=red!40!black,
	sharp corners=northwest,
	fonttitle=\sffamily\bfseries,
	title=~\thetcbcounter: #2,
	#1
}

% U.D.M
\newcommand{\amp}{\ensuremath{\mathrm{A}}}
\newcommand{\volt}{\ensuremath{\mathrm{V}}}
\newcommand{\meter}{\ensuremath{\mathrm{m}}}
\newcommand{\second}{\ensuremath{\mathrm{s}}}
\newcommand{\farad}{\ensuremath{\mathrm{F}}}
\newcommand{\henry}{\ensuremath{\mathrm{H}}}
\newcommand{\siemens}{\ensuremath{\mathrm{S}}}

% circuiti
\usepackage{circuitikz}
\usetikzlibrary{babel}

% disegni
\usepackage{pgfplots}
\pgfplotsset{width=10cm,compat=1.9}

% disponi codice
\usepackage{listings}
\usepackage[table]{xcolor}

\lstdefinestyle{codestyle}{
		backgroundcolor=\color{black!5}, 
		commentstyle=\color{codegreen},
		keywordstyle=\bfseries\color{magenta},
		numberstyle=\sffamily\tiny\color{black!60},
		stringstyle=\color{green!50!black},
		basicstyle=\ttfamily\footnotesize,
		breakatwhitespace=false,         
		breaklines=true,                 
		captionpos=b,                    
		keepspaces=true,                 
		numbers=left,                    
		numbersep=5pt,                  
		showspaces=false,                
		showstringspaces=false,
		showtabs=false,                  
		tabsize=2
}

\lstdefinestyle{shellstyle}{
		backgroundcolor=\color{black!5}, 
		basicstyle=\ttfamily\footnotesize\color{black}, 
		commentstyle=\color{black}, 
		keywordstyle=\color{black},
		numberstyle=\color{black!5},
		stringstyle=\color{black}, 
		showspaces=false,
		showstringspaces=false, 
		showtabs=false, 
		tabsize=2, 
		numbers=none, 
		breaklines=true
}

\lstdefinelanguage{javascript}{
	keywords={typeof, new, true, false, catch, function, return, null, catch, switch, var, if, in, while, do, else, case, break},
	keywordstyle=\color{blue}\bfseries,
	ndkeywords={class, export, boolean, throw, implements, import, this},
	ndkeywordstyle=\color{darkgray}\bfseries,
	identifierstyle=\color{black},
	sensitive=false,
	comment=[l]{//},
	morecomment=[s]{/*}{*/},
	commentstyle=\color{purple}\ttfamily,
	stringstyle=\color{red}\ttfamily,
	morestring=[b]',
	morestring=[b]"
}

% disponi sezioni
\usepackage{titlesec}

\titleformat{\section}
	{\sffamily\Large\bfseries} 
	{\thesection}{1em}{} 
\titleformat{\subsection}
	{\sffamily\large\bfseries}   
	{\thesubsection}{1em}{} 
\titleformat{\subsubsection}
	{\sffamily\normalsize\bfseries} 
	{\thesubsubsection}{1em}{}

% disponi alberi
\usepackage{forest}

\forestset{
	rectstyle/.style={
		for tree={rectangle,draw,font=\large\sffamily}
	},
	roundstyle/.style={
		for tree={circle,draw,font=\large}
	}
}

% disponi algoritmi
\usepackage{algorithm}
\usepackage{algorithmic}
\makeatletter
\renewcommand{\ALG@name}{Algoritmo}
\makeatother

% disponi numeri di pagina
\usepackage{fancyhdr}
\fancyhf{} 
\fancyfoot[L]{\sffamily{\thepage}}

\makeatletter
\fancyhead[L]{\raisebox{1ex}[0pt][0pt]{\sffamily{\@title \ \@date}}} 
\fancyhead[R]{\raisebox{1ex}[0pt][0pt]{\sffamily{\@author}}}
\makeatother

\begin{document}
% sezione (data)
\section{Lezione del 26-05-26}

% stili pagina
\thispagestyle{empty}
\pagestyle{fancy}

% testo
Nelle scorse lezioni abbiamo parlato dei filtri di trasmissione e ricezione nei sistemi di trasmissione in banda base ($G_T(f)$ e $G_R(f)$). In particolare, abbiamo detto che condizione per la massimizzazione dell'SNR è:
$$
G_R(f) = G_T(-f) = G_T^*(f)
$$
mentre per l'abbattimento dell'ISI deve essere:
$$
G(f) = G_T(f) G_R(f), \quad \left( g(t) = g_T(t) \circledast g_R(t) \right)
$$
con $G(f)$ approssimante l'impulso di Nyquist nullo a $\pm m$.

Quindi, scelto $g(t)$ ($G(f)$ impulso di Nyquist), ciò che facciamo è prendere:
$$
G(f) = G_T(f) G_R(f) = |G_T(f)|^2 \implies |G_T(f)| = \sqrt{G(f)}
$$
da cui deriva (a fase arbitraria):
$$
\begin{cases}
	G_T(f) = \sqrt{G(f)} e^{j \mathrm{arg}(G_T(f))} \\
	G_R(f) = \sqrt{G(f)} e^{-j \mathrm{arg}(G_T(f))} \\
\end{cases}
$$

Il filtro ricavato dall'impulso di Nyquist non è però valido, in quanto risulta che il caso più semplice dà:
$$
g_T(t) = \mathrm{rect} \left( \frac{t}{T_S} \right)
$$
da cui la densità septtrale di potenza:
$$
S_S(f) = \frac{\sigma_a^2}{T_S} |G_T(f) |^2
$$
va come la $\mathrm{sinc}^2$, cosa che è molto poco efficiente dal punto di vista dello spettro.

\subsubsection{Coseno rialzato}
Come abbiamo visto, la soluzione è da prendere $G(f)$ come l'interpolatore \textbf{coseno rialzato}:
$$
H(f)=
\begin{cases}
T, & |f|\le \frac{1-\beta}{2T} \\[6pt]
\frac{T}{2}\left[1+\cos\left(\frac{\pi T}{\beta}\left(|f|-\frac{1-\beta}{2T}\right)\right)\right], 
& \frac{1-\beta}{2T}<|f|\le \frac{1+\beta}{2T} \\[10pt]
0, & |f|>\frac{1+\beta}{2T}
\end{cases}
$$
dove $\beta \in [0, 1]$ viene detto fattore di rollof. Ricordiamo che il coseno rialzato inizia a decadere a $\frac{1 - \alpha}{2T}$, e finisce a $\frac{1 + \alpha}{2 T}$, con punto medio del \textit{rolloff} (o decadimento) in $\frac{1}{2T}$. La banda del filtro sarà $\frac{1 + \alpha}{2T}$, con valori di $\alpha$ intorno a $\sim 0.3$.

Notiamo che, scelto questo $G(f)$ (reale) si ricavano i filtri al trasmettitore e al ricevitore come \textit{radici di coseno rialzato}:
$$
G_R(f) = G_T(f) = \sqrt{G(f)}
$$
dall'andamento leggermente più "schiacciato" rispetto al semplice coseno rialzato.

Il coseno rialzato soddisfa la condizione di Nyquist (in quanto sovrapponendola a periodo $\frac{1}{T_S}$ dà una costante), e ha banda limitata controllabile dal fattore di rollof. 

In dominio tempo, il coseno rialzato dà:
$$
g(t) = \frac{\mathrm{sinc} \left( \frac{\pi t}{T_s} \right)}{\frac{\pi t}{T_s}} \cdot \frac{\cos\left( \alpha \pi \frac{t}{T_S} \right)}{1 - \left(\alpha 2 \frac{t}{T_S}\right)^2}
$$
cioè qualcosa che approssima sempre di più una $\mathrm{sinc}$ (a sinistra) per $\alpha \rightarrow 0$.

La scelta del parametro $\alpha$ è un compromesso fra la banda (che vogliamo più bassa possibile, cioè ridurre $\alpha$) e la robustezza alla frequenza di campionamento (che tende ad aumentrae $\alpha$): questo perchè la porzione di tempo per cui la $g(t)$ resta attorno a $1$ diminuisce al diminuire di $\alpha$, e questo rende più facile l'introduzione di errori da leggeri offset del campionatore a ricevitore.

\subsubsection{Calcolo del BER}
Veniamo al discutere il \textbf{BER} (\textit{Byte Error Rate}).
Questo sarà una misura della probabilità di errore sul singolo byte (o in genere sulla parola di codice). Spostiamoci quindi sul demodulatore, dove abbiamo in ingresso il segnale $x(k)$, riscalato a $z(k)$ col fattore $\sqrt{\beta} g(0)$.

La probabilità di errore sarà, dal teorema della probabilità totale:
$$
P(e) = \sum_{l=1}^M P(e | a_k = a^{(l)}) P( a_k = a^{(l)} )
$$
assumiamo una codifica PAM con costellazione a $\# M = 2$ (cioè con solo $2$ simboli). Vogliamo che la costellazione sia una mappa di \textit{Gray}, cioè che a parità di rumore minimizza la probabilità di errore (questo si ha equispaziando la costellazione). In questo caso si avrà:
$$
= \frac{1}{2} P(e | a_k = -1) + \frac{1}{2} P(e | a_k = 1)
$$
Con costellazione binaria, la decodifica del simbolo (cioè l'implementazione del demapper) è piuttosto triviale, in quanto:
$$
\begin{cases}
z(k) \geq 0 \implies \hat{a}_k = 1  \\
z(k) < 0 \implies \hat{a}_k = -1
\end{cases}
$$

Ora, a decidere i nostri errori sarà il rumore gaussiano, in quanto:
$$
z(k) = a_k + n(k), \quad n(k) \sim N\left( 0, \sigma_n'^2 \right)
$$
da cui:
$$
\begin{cases}
	(z(k) | a_k = 1) = 1 + n(k) \\
	(z(k) | a_k = -1) = -1 + n(k) \\
\end{cases}
$$
cioè i valori del segnale normalizzato (prima del demapper) al ricevitore, dato il valore trasmesso, sono rappresentati da distribuzioni gaussiane scostate verso sinistra o verso destra di modulo unitario.
L'errore $e$ si ha nei punti della campana che si trovano a sinistra dello $0$, per $a_k = 1$, e viceversa a destra dello $0$, per $a_k = -1$.

Quindi, visto che il rumore è gaussiano, il calcolo quantitativo della probabilità sara triviale in quanto vale:
$$
\begin{cases}
P(e | a_k = 1) = P(1 + n(k) < 0) = 1 - Q\left( \frac{0 - 1}{\sigma_n'} \right) \\
P(e | a_k = -1) = P(-1 + n(k) > 0) = Q\left( \frac{0 + 1}{\sigma_n'} \right) \\
\end{cases}
$$

Da questo ricaviamo la probabilità di errore totale come:
$$
P(e) = \frac{1}{2} P(e) + \frac{1}{2} P(e) = Q \left( \frac{1}{\sigma_n'} \right)
$$
Fra l'altro, la probabilità di errore non varia se si varia lo scostamento in ampiezza della costellazione (al massimo cambierà l'energia impiegata per la trasmissione).

\par\smallskip

Facciamo la stessa cosa per una codifica su 4 simboli, quindi con $Q = 2$ e $\# = 4$.
In questo caso ci rifaremo alla formula:
$$
P(e) = \sum_{l=1}^M P(e | a_k = a^{(l)}) P( a_k = a^{(l)} )
$$
sui valori a priori $a_k = -3, -1, 1, 3$. 

Cominciamo a calcolare le singole probabilità di errore a partire dal valore $a_k = 3$:
$$
P(e | a_k = 3) = 1 - Q \left( \frac{2 - 3}{\sigma_n'} \right) = Q \left( \frac{1}{\sigma_n'} \right) = P(e | a_k = -3)
$$
dove si ricava subito anche il caso simmetrico ($a_k = -3$).
Calcoliamo quindi la stessa cosa nel caso $a_k = 1$. In questo caso le cose sono leggermente diverse, in quanto abbiamo due zone in cui andiamo in errore (sotto $0$ e sopra $2$):
$$
P(e | a_k = 1) = Q\left( \frac{0 - 1}{\sigma_n'} \right) + 1 - Q\left( \frac{2 - 1}{\sigma_n'} \right) = 2 Q\left( \frac{1}{\sigma_n'} \right) = P(e | a_k = -1)
$$
dove si ricava subito anche il caso simmetrico ($a_k = -1$).

A questo punto la probabilità di errore sulla parola di codice è data da:
$$
P(e) = \frac{1}{4} \left( 2 Q \left( \frac{1}{\sigma_n'} \right) + 4 Q \left( \frac{1}{\sigma_n'} \right) \right) = \frac{3}{2} Q \left( \frac{1}{\sigma_n'} \right)
$$

Notiamo quindi che non è detto che la probabilità di errore sul bit corrisponda alla probabilità di errore sul simbolo. In particolare, vale la stima:
$$
\frac{P(e)}{\log_2 (\# M)} \leq P_{bit}(e) \leq P(e)
$$

In generale, quindi, vale che le probabilità di errore BER per le codifiche PAM vanno come:
$$
P(e) = k Q\left( \frac{1}{\sigma_n'} \right)
$$
con $k$ dipendente dalla dimensione dell'alfabeto di simboli ($M$) usato.
A questo punto resta da documentare come si ricava la varianza $\sigma_n^2$, da cui si ricava la deviazione a denominatore. Questa era già stata calcolata come:
$$
\sigma_n'^2 = \frac{\sigma_n^2}{\beta g^2(0)}, \quad \sigma_n^2 = \frac{N_0}{2} E_{g_R}
$$
con $N_0 = k_B T$ dal rumore termico.
Questo dipende dall'energia dell'interpolatore a ricevitore, che sarà:
$$
E_{g_R} = \int_{-\infty}^{\infty} | G_R(f) |^2 \, df
$$

Prendendo di usare l'interpolatore a coseno rialzato, si ha:
$$
\sigma_n'^2 = \frac{N_0}{2 \beta}
$$

\subsubsection{Considerazioni di potenza}

Calcoliamo allora l'energia di trasmissione di un simbolo con interpolatore a coseno rialzato, e filtri adattati. Questa sarà:
$$
E_s = P_S T_S = T_S \int_{-\infty}^{\infty} S_S(f) \, df = \frac{M^2 - 1}{3} \int_{-\infty}^{\infty} | G_T(f) |^2 \, df = \frac{M^2 - 1}{3}
$$
nel caso del $G(f)$ coseno rialzato. La densità spettrale di potenza sarà:
$$
S_S(f) = \frac{1}{T_S} \sigma_n^2 | G_T(f) |^2 = \frac{1}{T_S} \frac{M^2 - 1}{3} | G_T(f) |^2
$$

L'energia del simbolo ricevuto sarà quella del simbolo trasmesso, col fattore di attenuazione:
$$
E_r = \beta E_s
$$
da cui il rapport:
$$
\frac{E_r}{N_0} = \beta \frac{M^2 - 1}{3} \frac{1}{N_0} \implies N_0 = \beta \frac{M^2 - 1}{3} \frac{1}{\left(\frac{E_r}{N_0}\right)}
$$
Sostituendo in $\sigma_n'^2$, si ottiene:
$$
\sigma_n'^2 = \frac{N_0}{2 \beta} = \frac{1}{2} \frac{M^2 - 1}{3} \frac{1}{\left(\frac{E_r}{N_0}\right)}
$$

Abbiamo tra l'altro un indicazione in funzione della $M$:
$$
P(e) =  Q \left( \frac{1}{\sigma_n'} \right) = Q \left( \sqrt{\frac{6 \frac{E_r}{N_0} }{M^2 - 1}} \right) 
$$

Facciamo un esempio: presi due sistemi, uno con $\# M = 2$ e l'altro con $\# M = 4$ (un sistema binario e un sistema quaternario), si ha:
$$
\begin{cases}
	M = 2 \implies P(e) \sim Q\left( \sqrt{2 \frac{E_r}{N_0}} \right) \\
	M = 4 \implies P(e) \sim Q\left( \sqrt{\frac{2}{5} \frac{E_r}{N_0}} \right)
\end{cases}
$$

Per cui se vogliamo sviluppare 2 sistemi, uno binario e uno quaternario, con la stessa probabilità di errore, dobbiamo imporre:
$$
P_2 (e) = P_4 (e) \implies Q\left( \sqrt{2 \frac{E_r}{N_0}} \right) = Q\left( \sqrt{\frac{2}{5} \frac{E_r}{N_0}} \right)
$$
da cui si ricava riguardo alle energie:
$$
E_r \Big|_{\# M = 4} = 5 E_r \Big|_{\# M = 2}
$$
cioè per avere la stessa probabilità di errore bisogna aumentare di un fattore di $5$ l'energia impiegata nel sistema quaternario.

Infine, ricordiamo riguardo alla banda:
$$
B_s = \frac{R_d}{r \log_2 (\# M)}
$$
cioè all'aumentare di $\# M$ si riduce la banda necessaria alla trasmisione dello stesso segnale. Questo compromesso (fra banda e probabilità di errore) è lo stesso a cui accennavamo nelle scorse lezioni. 


\end{document}
